\chapter{晶体的下降性质}
\markboth{第 3 章\ 晶体的下降性质}{第 3 章\ 晶体的下降性质}

本节中，所有预栈（包括 DG 概形与 DG ind-概形）除明确说明外都假定为局部几乎有限型。

本节的目标是建立一系列关于晶体在 DG 概形与 DG ind-概形上行为的性质。这些性质包括：用无穷小群胚解释（左、右）晶体；h-下降；晶体范畴的单子描述；右晶体与左晶体的归纳函子。

\section{无穷小群胚}

本小节中，$X$ 是局部几乎有限型的 DG ind-概形，见 [GR1, Sect.~1.7.1]。

\ob{3.1.1}{n:3.1.1}考虑单纯预栈 $X_\bullet/X_{\dR}$，即映射 $X\to X_{\dR}$ 的 Čech 神经。前面已经指出，每个 $X^i/X_{\dR}$ 都是 $X^i$ 沿主对角线的形式完备化。特别地，所有 $X^i/X_{\dR}$ 也都属于 $\DGIndSch_{\mathrm{laft}}$。

我们把
\[
X\times_{X_{\dR}}X\rightrightarrows X
\]
称为 $X$ 的无穷小群胚。

\ob{3.1.2}{n:3.1.2}考虑余单纯范畴 $\IndCoh(X_\bullet/X_{\dR})$。

\begin{othm}{命题}{3.1.3}{n:3.1.3}
由增广定义的函子
\[
\Crys^r(X)\to\Tot\big(\IndCoh(X_\bullet/X_{\dR})\big)
\]
是等价。
\end{othm}

\begin{othm}{注}{3.1.4}{n:3.1.4}
注意：由引理 2.3.11，把 $X$ 换成任意经典形式光滑的对象 $Y\in\PreStk_{\mathrm{laft}}$，命题的断言仍成立。
\end{othm}

\begin{proof}
只需证明：对任意 $S\in\DGSch_{\mathrm{aff}}^{\mathrm{ft}}$ 及映射 $S\to X_{\dR}$，函子
\[
\IndCoh(S)\to\Tot\Big(\IndCoh\big(S\times_{X_{\dR}}X^\bullet/X_{\dR}\big)\Big)
\]
是等价。

注意，单纯预栈 $S\times_{X_{\dR}}(X_\bullet/X_{\dR})$ 是映射
\begin{equation}
S\times_{X_{\dR}}X\to S\tag{3.1}
\end{equation}
的 Čech 神经。

注意 $S\times_{X_{\dR}}X$ 等同于 $S\times X$ 沿映射 $\clred S\to S\times X$ 的形式完备化，其中 $\clred S\to X$ 是对应于 $S\to X_{\dR}$ 的映射。特别地，我们得到：(3.1) 中的映射是 ind-固有的（ind-固有性的概念见 [GR1, Sect.~2.7.4]）且满。

于是，断言由 [GR1, Lemma 2.10.3] 推出。
\end{proof}

\section{晶体的 fppf 下降与 h-下降}

\ob{3.2.1}{n:3.2.1}回忆范畴 $\DGSch_{\mathrm{aff}}^{\mathrm{aft}}$ 上的 h-拓扑 [IndCoh, Sect.~8.2]。它由 Zariski 覆盖与固有满覆盖生成。

考虑函子
\[
\Crys^r_{\DGSch_{\mathrm{aff}}^{\mathrm{aft}}}:=\Crys^r_{\PreStk_{\mathrm{laft}}}|_{\DGSch_{\mathrm{aff}}^{\mathrm{aft}}}:(\DGSch_{\mathrm{aff}}^{\mathrm{aft}})^{\mathrm{op}}\to\DGCat.
\]
我们将证明：

\begin{othm}{命题}{3.2.2}{n:3.2.2}
函子 $\Crys^r_{\DGSch_{\mathrm{aff}}^{\mathrm{aft}}}$ 满足 h-下降。
\end{othm}

\begin{proof}
我们将证明 $\Crys^r_{\DGSch_{\mathrm{aff}}^{\mathrm{aft}}}$ 满足平展下降与固有满下降。

平展下降是清楚的：若 $S'\to S$ 是 $\DGSch_{\mathrm{aff}}^{\mathrm{aft}}$ 中的平展覆盖，则相应的映射 $S'_{\dR}\to S_{\dR}$ 是 $\PreStk_{\mathrm{laft}}$ 中概形的、平展的、满的映射。特别地，它是 fppf 拓扑的覆盖，于是断言由 $\IndCoh$ 的 fppf 下降推出，见 [IndCoh, Corollary 10.4.5]。

于是，设 $S'\to S$ 是固有满映射。考虑 $\PreStk_{\mathrm{laft}}$ 中的双单纯对象
\[
(S'^\bullet/S)_\star/(S'^{\bullet}_{\dR}/S_{\dR}),
\]
即 $S'\to S$ 的 Čech 神经的逐项无穷小群胚。具体地说，它是这样的双单纯对象：其 $(p,q)$-单形由映射
\[
S'^p/S\to S'^p_{\dR}/S_{\dR}
\]
的 Čech 神经的 $q$-单形给出；因此 $\star$ 标记指标 $q$，$\bullet$ 标记指标 $p$。

由命题 3.1.3，只需证明复合函子
\begin{equation}
\Crys^r(S):=\IndCoh(S_{\dR})\to\Tot\big(\IndCoh((S'^\bullet/S)_{\dR})\big)\to
\Tot\big(\IndCoh((S'^\bullet/S)_\star/(S'^\bullet/S)_{\dR})\big)\tag{3.2}
\end{equation}
是等价。

然而注意，$\PreStk_{\mathrm{laft}}$ 的双单纯对象有典范同构
\[
(S'^\bullet/S)_\star/(S'^{\bullet}_{\dR}/S_{\dR})\simeq(S'^\star/S_{\dR})_\bullet/(S'^\star/S_{\dR}),
\]
后者是余单纯对象
\[
(S'^\star/S_{\dR})\to(S^\star/S_{\dR})
\]
的逐项 Čech 神经。

(3.2) 中的映射可以改写为
\[
\Crys^r(S):=\IndCoh(S_{\dR})\to\Tot\big(\IndCoh(S^\bullet/S_{\dR})\big)\to\Tot\big(\IndCoh((S'^\star/S_{\dR})_\bullet/(S^\star/S_{\dR}))\big).
\]
再次应用命题 3.1.3，我们得到：只需证明对每个 $i$，映射
\[
\IndCoh(S^i/S_{\dR})\to\Tot\Big(\IndCoh\big((S'^i/S'^i_{\dR})_\bullet/(S^i/S_{\dR})\big)\Big)
\]
是等价。

然而注意，映射
\[
S'^i/S_{\dR}\to S^i/S_{\dR}
\]
是 ind-固有的且满。于是断言由 [GR1, Lemma 2.10.3] 推出。
\end{proof}

\ob{3.2.3}{n:3.2.3}考虑范畴 $\DGSch_{\mathrm{aff}}^{\mathrm{aft}}$ 上的 fppf 拓扑，它由 $\DGSch_{\mathrm{aff}}$ 上的 fppf 拓扑诱导（见 [GL:Stacks, Sect.~2.2]）。注意每个 fppf 覆盖特别是 h-覆盖。因此得到：

\begin{othm}{推论}{3.2.4}{n:3.2.4}
函子 $\Crys^r_{\DGSch_{\mathrm{aff}}^{\mathrm{aft}}}$ 满足 fppf 下降。
\end{othm}

如同 [IndCoh, Theorem 8.3.2]，fppf 下降是 Nisnevich 下降与有限平坦下降的结合\footnote{这个观察由 J. Lurie 向我们解释。}。特别地，我们在命题 3.2.2 的证明中建立 fppf 下降时，并没有用到「每个 fppf 覆盖也是 h-覆盖」这一事实。

\ob{3.2.5}{n:3.2.5}$\DGSch_{\mathrm{aff}}^{\mathrm{aft}}$ 上的 fppf（相应地，h-）拓扑诱导全子范畴
\[
{}^{<\infty}\DGSch_{\mathrm{aff}}^{\mathrm{aft}}\subset\DGSch_{\mathrm{aff}}^{\mathrm{aft}}
\]
上的 fppf（相应地，h-）拓扑。

命题 3.2.2 蕴含：

\begin{othm}{推论}{3.2.6}{n:3.2.6}
函子
\[
\Crys^r_{{}^{<\infty}\DGSch_{\mathrm{aff}}^{\mathrm{aft}}}:=\Crys^r_{\PreStk_{\mathrm{laft}}}|_{{}^{<\infty}\DGSch_{\mathrm{aff}}^{\mathrm{aft}}}:({}^{<\infty}\DGSch_{\mathrm{aff}}^{\mathrm{aft}})^{\mathrm{op}}\to\DGCat
\]
在 ${}^{<\infty}\DGSch_{\mathrm{aff}}^{\mathrm{aft}}$ 上满足 h-下降，特别地满足 fppf 下降。
\end{othm}

于是由 [Lu0, Corollary 6.2.3.5] 得到：

\begin{othm}{推论}{3.2.7}{n:3.2.7}
设 $Y_1\to Y_2$ 是 $\PreStk_{\mathrm{laft}}$ 中的映射，且在 h-拓扑下是满射。则自然映射
\[
\Crys(Y_2)\to\Tot\big(\Crys(Y^1_\bullet/Y_2)\big)
\]
是等价。
\end{othm}

\section{右晶体的归纳函子}

\ob{3.3.1}{n:3.3.1}记 $p_s,p_t$ 为两个投影
\[
X\times_{X_{\dR}}X\rightrightarrows X.
\]
注意映射 $p_i$（$i=s,t$）是 ind-固有的。因此函子 $p_i^!$ 有左伴随 $(p_i)_*^{\IndCoh}$，见 [GR1, Corollary 2.8.3]。

\begin{othm}{命题}{3.3.2}{n:3.3.2}
(a) 遗忘函子
\[
\oblv^r_X:\Crys(X)\to\IndCoh(X)
\]
有左伴随，记为 $\indf^r_X$。

(b) 有函子的典范同构
\[
\oblv^r_X\circ\indf^r_X\simeq(p_t)_*^{\IndCoh}\circ(p_s)^!.
\]

(c) 伴随对
\[
\indf^r_X:\IndCoh(X)\rightleftarrows\Crys^r(X):\oblv^r_X
\]
是单子的；也就是说，从 $\Crys^r(X)$ 到 $\IndCoh(X)$ 上关于单子 $\oblv^r_X\circ\indf^r_X$ 的模范畴的自然函子是等价。
\end{othm}

\begin{proof}
由命题 3.1.3 与 [Lu2, Theorem 6.2.4.2]，只需证明余单纯范畴
\[
\IndCoh(X_\bullet/X_{\dR})
\]
满足 Beck--Chevalley 条件，即：对每个 $n$，余面映射
\[
d^0:\IndCoh(X^n/X_{\dR})\to\IndCoh(X^{n+1}/X_{\dR})
\]
有左伴随，记为 $t^0$；并且对 $\Delta$ 中每个映射 $[m]\to[n]$，图表
\[
\begin{tikzcd}[column sep=large, row sep=large]
 \IndCoh(X^m/X_{\dR}) & \IndCoh(X^{m+1}/X_{\dR}) \arrow[l, "t^0"'] \\
 \IndCoh(X^n/X_{\dR}) \arrow[u] & \IndCoh(X^{n+1}/X_{\dR}) \arrow[l, "t^0"'] \arrow[u]
\end{tikzcd}
\]
（它先验地只交换到自然变换）实际上交换。

在这一情形，Beck--Chevalley 条件归结为 DG ind-概形间 ind-固有态射的 $*$-推前与 $!$-拉回之间的伴随性与基变换，由 [GR1, Proposition 2.9.2] 给出。
\end{proof}

\begin{othm}{推论}{3.3.3}{n:3.3.3}
范畴 $\Crys^r(X)$ 是紧生成的。
\end{othm}

\begin{proof}
紧生成元的集合由把 $\indf^r_X$ 作用于 $\IndCoh(X)$ 的紧生成元 $\Coh(X)$ 得到（见 [GR1, Corollary 2.4.4]）。
\end{proof}

\section{左晶体的归纳函子与无穷小群胚}

\ob{3.4.1}{n:3.4.1}由引理 1.2.4，对光滑经典概形 $X$，命题 3.1.3 的类比对左晶体成立，即函子
\begin{equation}
\Crys^l(X)=\QCoh(X_{\dR})\to\Tot\big(\QCoh(X^\bullet/X_{\dR})\big)\tag{3.3}
\end{equation}
是等价。

由命题 3.1.3，右晶体的类似断言对任意 DG 概形（乃至 DG ind-概形）$X$ 成立。然而，左晶体并非如此。

\ob{3.4.2}{n:3.4.2}我们断言：

\begin{othm}{命题}{3.4.3}{n:3.4.3}
若 DG 概形 $X$ 最终余连接，则函子 (3.3) 是等价。
\end{othm}

\begin{othm}{注}{3.4.4}{n:3.4.4}
可以证明，命题的断言对任意局部几乎有限型的 DG 概形 $X$ 都成立。但证明更为复杂。此外，由引理 2.1.7，该断言对任意经典形式光滑的预栈也成立。
\end{othm}

\begin{othm}{例}{3.4.5}{n:3.4.5}
考虑 DG 概形 $X=\Spec(k[\alpha])$，其中 $\alpha$ 位于次数 $-2$。这是构造涉及 $\Crys^l(X)$ 的断言的反例时值得记住的好例子。
\end{othm}

\begin{annote}
这个例子值得记住的原因：$\alpha$ 落在次数 $-2$，使 $k[\alpha]$ 在负次数有无界同伦群，从而 $X$ 不最终余连接。直观上，此时 $X$ 沿自身主对角线的形式完备化不再是 $X$ 本身，$\Tot(\QCoh(X^\bullet/X_{\dR}))$ 会比 $\Crys^l(X)$ 多出对象；这正是命题 3.4.3 需要最终余连接假设的原因。\hfill【译者按（上下文推断，无外部依据）】
\end{annote}


\begin{proof}[命题 3.4.3 的证明]
有函子的交换图
\[
\begin{tikzcd}[column sep=large, row sep=large]
 \Crys^l(X) \arrow[r] \arrow[d, "\Upsilon_X"'] & \Tot\big(\QCoh(X^\bullet/X_{\dR})\big) \arrow[d, "\Tot(\Upsilon_{X^\bullet/X_{\dR}})"] \\
 \Crys^r(X) \arrow[r] & \Tot\big(\IndCoh(X^\bullet/X_{\dR})\big),
\end{tikzcd}
\]
其中左竖直映射与下水平映射是等价。于是得到：$\Crys^l(X)$ 是 $\Tot(\QCoh(X^\bullet/X_{\dR}))$ 的收缩。

回忆：若 $Z$ 是最终余连接的 DG 概形，则函子
\[
\Upsilon_Z:\QCoh(Z)\to\IndCoh(Z)
\]
是全忠实的（见 [IndCoh, Corollary 9.6.3]）。于是由 [GR1, Propositions 7.1.3 and 7.4.5]，对最终余连接 DG 概形沿 Zariski 闭子集的形式完备化，同样的结论也成立。因此函子
\[
\Upsilon_{X^i/X_{\dR}}:\QCoh(X^i/X_{\dR})\to\IndCoh(X^i/X_{\dR})
\]
全忠实。于是上述交换图中的函子 $\Tot(\Upsilon_{X^\bullet/X_{\dR}})$ 也全忠实。但恒等函子是它的收缩，故它本质满。
\end{proof}

\ob{3.4.6}{n:3.4.6}对 DG 概形 $X$，定义函子
\[
\indf^l_X:\QCoh(X)\to\Crys^l(X)
\]
为
\[
\indf^l_X:=(\Upsilon_{X_{\dR}})^{-1}\circ\indf^r_X\circ\Upsilon_X.
\]
我们断言：

\begin{othm}{引理}{3.4.7}{n:3.4.7}
若 $X$ 是最终余连接的 DG 概形，则函子 $(\indf^l_X,\oblv^l_X)$ 互为伴随。
\end{othm}

\begin{othm}{注}{3.4.8}{n:3.4.8}
若去掉 $X$ 最终余连接的假设，引理的断言不成立。事实上，此时函子 $\indf^l_X$ 不保持紧性。
\end{othm}

\begin{proof}[引理 3.4.7 的证明]
回忆（见 [IndCoh, Sect.~9.6.6]）：对最终余连接的 $X$，函子 $\Upsilon_X$ 有右伴随，记为 $\Xi^{\vee}_X$；此外，$\Upsilon_X$ 本身全忠实。

由此得到 $\indf^l_X$ 的右伴随为
\[
\Xi^{\vee}_X\circ\oblv^r_X\circ\Upsilon_{X_{\dR}}\simeq\Xi^{\vee}_X\circ\Upsilon_X\circ\oblv^l_X\simeq\oblv^l_X,
\]
这正是所要的。
\end{proof}

在引理 3.4.7 的证明过程中，我们还看到：

\begin{othm}{引理}{3.4.9}{n:3.4.9}
函子 $\oblv^l_X$ 典范同构于
\[
\Xi^{\vee}_X\circ\oblv^r_X\circ\Upsilon_{X_{\dR}}.
\]
\end{othm}

\ob{3.4.10}{n:3.4.10}现在断言：

\begin{othm}{命题}{3.4.11}{n:3.4.11}
设 $X$ 是最终余连接的 DG 概形。则伴随对
\[
\indf^l_X:\QCoh(X)\rightleftarrows\Crys^l(X):\oblv^l_X
\]
是单子的；也就是说，从 $\Crys^l(X)$ 到 $\QCoh(X)$ 上关于单子 $\oblv^l_X\circ\indf^l_X$ 的模范畴的自然函子是等价。
\end{othm}

\begin{proof}
需要说明 Barr--Beck--Lurie 定理的条件成立。函子 $\oblv^l_X$ 连续，因此与所有余极限交换。$\oblv^l_X$ 保守这一事实由引理 2.2.6 给出。
\end{proof}
